\documentclass[11pt,a4paper]{article}
\usepackage[margin=0.95in]{geometry}
\usepackage[T1]{fontenc}
\usepackage{lmodern}
\usepackage{amsmath,amssymb}
\usepackage{booktabs}
\usepackage{xcolor}
\usepackage{fancyhdr}
\usepackage{hyperref}
\setlength{\parskip}{0.46em}
\setlength{\parindent}{0pt}
\setlength{\headheight}{15pt}
\linespread{1.07}
\pagestyle{fancy}
\fancyhf{}
\lhead{CST444 Soft Computing - Module 4}
\rhead{Exam-Ready Beginner Guide}
\cfoot{\thepage}
\hypersetup{colorlinks=true,linkcolor=blue,urlcolor=blue}

\newcommand{\examtip}[1]{%
\vspace{0.2em}
\noindent\fcolorbox{black}{gray!12}{\parbox{0.97\linewidth}{\textbf{Exam Tip:} #1}}%
\vspace{0.2em}
}

\begin{document}

\section*{Module 4 Topic Map from OCR Questions}
\textbf{Frequently repeated concepts}
\begin{itemize}
\item Fuzzy Inference System (FIS): architecture, operation, Mamdani and Sugeno models.
\item Design steps of fuzzy logic controller (FLC).
\item Genetic Algorithm (GA): flowchart, operators, and stopping conditions.
\item Selection operators in GA.
\item Crossover operators: single-point, two-point, uniform, three-parent, shuffle, precedence-preservative.
\item Encoding methods: binary, value/real, permutation, tree.
\end{itemize}

\textbf{High-mark asks}
\begin{itemize}
\item Explain Mamdani and Sugeno FIS with diagram and comparison.
\item Explain selection and crossover operators with examples.
\item Explain encoding schemes and stopping criteria in GA.
\item Design-type question: Mamdani controller for washing machine.
\end{itemize}

\textbf{Must-study-first order}
\begin{enumerate}
\item FIS blocks and flow.
\item Mamdani vs Sugeno.
\item GA pipeline from population to next generation.
\item Selection and crossover.
\item Mutation and stopping criteria.
\item Encoding schemes.
\end{enumerate}

\section*{Core Theory (Teaching-First, Exam-Formal)}

\subsection*{1) Fuzzy Inference System (FIS)}
FIS is a rule-based mapping from input to output using fuzzy logic. It uses IF-THEN rules and operators like AND/OR.

\textbf{Main blocks in FIS/FLC}
\begin{itemize}
\item Fuzzifier
\item Rule base
\item Database/knowledge base
\item Decision-making unit (inference engine)
\item Defuzzifier
\end{itemize}

\begin{center}
\framebox[0.97\linewidth][c]{
\begin{minipage}{0.93\linewidth}
\vspace{0.08cm}
\centerline{\textbf{FIS Block Diagram}}
\centerline{Crisp Input $\rightarrow$ Fuzzifier $\rightarrow$ Inference Engine + Rule Base}
\centerline{$\rightarrow$ Aggregation $\rightarrow$ Defuzzifier $\rightarrow$ Crisp Output}
\vspace{0.08cm}
\end{minipage}}
\end{center}

\subsection*{2) Mamdani vs Sugeno FIS}
\textbf{Mamdani:} consequent is fuzzy set. Defuzzification is mandatory.

\textbf{Sugeno:} consequent is linear/constant function of inputs. Computationally compact and optimization-friendly.

\textbf{Typical Mamdani steps}
\begin{enumerate}
\item Fuzzify inputs.
\item Evaluate rule strengths.
\item Clip/scale output membership functions.
\item Aggregate all rule outputs.
\item Defuzzify to crisp output.
\end{enumerate}

\subsection*{3) FLC Design Steps}
\begin{enumerate}
\item Identify input, output, and state variables.
\item Partition universe and assign linguistic labels.
\item Select membership functions.
\item Build rule base.
\item Normalize/scaling choices.
\item Fuzzification.
\item Inference and aggregation.
\item Defuzzification.
\end{enumerate}

\subsection*{4) Genetic Algorithm Basics}
\textbf{GA cycle:} initialize population $\rightarrow$ fitness evaluation $\rightarrow$ selection $\rightarrow$ crossover $\rightarrow$ mutation $\rightarrow$ next generation.

\begin{center}
\framebox[0.97\linewidth][c]{
\begin{minipage}{0.93\linewidth}
\vspace{0.08cm}
\centerline{\textbf{GA Flow}}
\centerline{Initialize Population $\rightarrow$ Evaluate Fitness $\rightarrow$ Select Parents}
\centerline{$\rightarrow$ Crossover $\rightarrow$ Mutation $\rightarrow$ New Population $\rightarrow$ Stop?}
\vspace{0.08cm}
\end{minipage}}
\end{center}

\subsection*{5) Selection, Crossover, Mutation, Encoding}
\textbf{Selection methods (as in notes):}
Roulette-wheel, rank selection, tournament selection, stochastic universal sampling.

\textbf{Crossover methods:}
Single-point, two-point, multi-point, uniform, three-parent, shuffle, precedence-preservative.

\textbf{Mutation examples:}
Bit-flip, interchanging (swap), reversing.

\textbf{Encoding types:}
Binary, value/real, permutation, tree encoding.

\subsection*{6) Stopping Criteria in GA}
\begin{itemize}
\item Maximum generation reached.
\item Maximum elapsed time reached.
\item No significant fitness improvement.
\item Reaching desired objective/fitness threshold.
\item Saturation in best-individual behavior.
\end{itemize}

\examtip{For GA answers, always write operators in this order: selection, crossover, mutation. This matches standard marking flow.}

\section*{Numericals and Derivations (Stepwise)}

\subsection*{N1) Roulette-wheel selection probability (worked pattern)}
\textbf{Given:} Fitness values for four chromosomes: $[2,6,3,1]$.

\textbf{Find:} Selection probabilities.

\textbf{Formula:}
\[
P_i=\frac{f_i}{\sum f_i}
\]

\textbf{Substitution:}
\[
\sum f_i=2+6+3+1=12
\]
\[
P=[2/12,6/12,3/12,1/12]=[0.1667,0.5,0.25,0.0833]
\]

\textbf{Final answer:} Probabilities are $16.67\%$, $50\%$, $25\%$, $8.33\%$.

\subsection*{N2) Single-point crossover example}
\textbf{Given:} Parent1 = 11001010, Parent2 = 00110111, cut after 4th bit.

\textbf{Step}
\[
\text{Child1}=1100|0111=11000111,
\quad
\text{Child2}=0011|1010=00111010
\]

\textbf{Final answer:} Children are 11000111 and 00111010.

\subsection*{N3) Triangle inference pattern (I, R, IR, T) for angles $120^\circ,50^\circ,10^\circ$}
\textbf{Given:} $X=120, Y=50, Z=10$ with $X\ge Y\ge Z$, $X+Y+Z=180$.

\textbf{Find:} Membership values for isosceles (I), right (R), isosceles-right (IR), and other triangle (T).

\textbf{Method (notes style):}
Use the given membership formulas for each fuzzy shape and substitute angle values; then combine using min/max where required (for IR use logical intersection).

\textbf{Exam-ready result format:}
\begin{itemize}
\item Compute $\mu_I$ from closeness of $Y$ and $Z$.
\item Compute $\mu_R$ from closeness of one angle to $90^\circ$.
\item Compute $\mu_{IR}=\min(\mu_I,\mu_R)$.
\item Compute residual/other-triangle membership from complement logic.
\end{itemize}

\section*{Solved Questions for Module 4 (Self-Contained)}

\subsection*{Part A / 3-Mark Questions (Each Answer Has 6 Scoring Points)}

\textbf{Q1. Draw flow chart and explain steps of Genetic Algorithm.}
\begin{enumerate}
\item Start with random initial population of chromosomes.
\item Evaluate fitness of each chromosome.
\item Select parents based on fitness using a selection method.
\item Apply crossover to create offspring.
\item Apply mutation to maintain diversity.
\item Form next generation and repeat until stopping criterion.
\end{enumerate}

\textbf{Q2. Explain any 3 mutation techniques with examples.}
\begin{enumerate}
\item Bit-flip mutation: 0 changes to 1 or 1 to 0 at selected locus.
\item Interchange mutation: swap two positions in chromosome.
\item Reversing mutation: reverse order of genes in selected segment.
\item Mutation helps recover lost genetic material.
\item Mutation prevents premature convergence.
\item Mutation rate is kept low to preserve good schemata.
\end{enumerate}

\textbf{Q3. Explain stopping conditions for GA flow.}
\begin{enumerate}
\item Stop if maximum generation count is reached.
\item Stop if maximum execution time is reached.
\item Stop if desired fitness/target objective is achieved.
\item Stop if best fitness does not improve for many generations.
\item Stop if population diversity collapses and search saturates.
\item Practical systems may combine multiple stop conditions.
\end{enumerate}

\textbf{Q4. What is Fuzzy Inference System (FIS)?}
\begin{enumerate}
\item FIS is a fuzzy rule-based decision mapping system.
\item It uses IF-THEN rules with linguistic variables.
\item Inputs may be crisp or fuzzy.
\item Fuzzifier converts crisp inputs to fuzzy values.
\item Inference engine applies rule base and combines outputs.
\item Defuzzifier converts final fuzzy output to crisp decision.
\end{enumerate}

\textbf{Q5. Differentiate value encoding and permutation encoding.}
\begin{enumerate}
\item Value encoding stores genes as real/float values.
\item Permutation encoding stores genes as ordered sequence.
\item Value encoding is used in continuous parameter optimization.
\item Permutation encoding is used in scheduling/order problems.
\item Crossover/mutation operators differ for these encodings.
\item Permutation operators must preserve sequence validity.
\end{enumerate}

\textbf{Q6. Explain single-point and two-point crossover briefly.}
\begin{enumerate}
\item Single-point crossover uses one cut location.
\item Tails after cut are exchanged between parents.
\item Two-point crossover uses two cut locations.
\item Middle segment between cuts is exchanged.
\item Two-point crossover generally explores more combinations.
\item Both are recombination operators for offspring generation.
\end{enumerate}

\subsection*{Part B / 14-Mark Questions (Each Answer Has 14 Scoring Points)}

\textbf{Q1. What is FIS? Illustrate Mamdani and Sugeno FIS with examples.}
\begin{enumerate}
\item FIS is a fuzzy rule-based mapping from input space to output space.
\item It consists of fuzzifier, rule base, database, inference engine, and defuzzifier.
\item Mamdani model uses fuzzy sets in rule consequents.
\item Sugeno model uses linear/constant output function in consequents.
\item Mamdani is intuitive and widely accepted in expert-style control.
\item Sugeno is computationally efficient and optimization-friendly.
\item In Mamdani, output membership functions are aggregated.
\item Mamdani requires explicit defuzzification at output stage.
\item In Sugeno zero-order model, consequent is a constant value.
\item In Sugeno first-order model, consequent is linear in inputs.
\item Mamdani example rule: IF temperature is high THEN fan speed is fast (fuzzy).
\item Sugeno example rule: IF temperature is high THEN speed = $aT+b$.
\item Mamdani suits linguistic interpretability; Sugeno suits adaptive tuning.
\item Conclusion: both share fuzzy front-end but differ in consequent modeling and output computation.
\end{enumerate}

\textbf{Q2. Explain any three selection techniques in GA.}
\begin{enumerate}
\item Selection chooses parents for reproduction based on fitness.
\item Roulette-wheel gives probability proportional to fitness.
\item Rank selection assigns probabilities based on ranking, not raw fitness.
\item Tournament selection selects best individual among random subset.
\item Stochastic universal sampling reduces selection variance.
\item Strong selection speeds convergence but can reduce diversity.
\item Weak selection preserves diversity but slows convergence.
\item Roulette-wheel is simple and popular in basic GA.
\item Rank method helps avoid domination by very fit individuals.
\item Tournament method is efficient and easy to implement.
\item Selection pressure controls exploration vs exploitation balance.
\item Parent pool created from selected individuals feeds crossover stage.
\item Selection does not create new genes by itself.
\item Conclusion: selection strategy strongly affects convergence speed and robustness.
\end{enumerate}

\textbf{Q3. Explain uniform, three-parent, shuffle and precedence-preservative crossover with examples.}
\begin{enumerate}
\item Uniform crossover decides gene-by-gene parent source.
\item It uses random mask or probability at each locus.
\item Three-parent crossover combines information from three parents.
\item Third parent helps resolve conflicts when first two differ.
\item Shuffle crossover shuffles gene order before crossover.
\item Then normal crossover is applied and order is unshuffled.
\item This reduces positional bias of fixed crossover points.
\item Precedence-preservative crossover is used for ordering/scheduling tasks.
\item It preserves relative precedence constraints.
\item It avoids invalid offspring in permutation problems.
\item These methods are chosen based on chromosome representation.
\item Binary/value problems often use uniform and multipoint methods.
\item Permutation problems prefer precedence-preservative variants.
\item Conclusion: crossover design must match encoding to preserve feasibility and diversity.
\end{enumerate}

\textbf{Q4. Explain various encoding schemes used in GA.}
\begin{enumerate}
\item Encoding maps candidate solution to chromosome form.
\item Binary encoding uses bit strings (0/1 genes).
\item Value/real encoding uses floating-point values.
\item Permutation encoding uses ordered lists.
\item Tree encoding represents expressions/program structures.
\item Binary encoding is easy for classical genetic operators.
\item Real encoding is effective for continuous optimization.
\item Permutation encoding suits route/schedule/order tasks.
\item Tree encoding is common in genetic programming.
\item Operator design depends on encoding validity constraints.
\item Mutation/crossover for permutation must avoid duplicates.
\item Real encoding often uses arithmetic crossover and perturbation mutation.
\item Good encoding improves search efficiency and feasibility handling.
\item Conclusion: problem type determines best encoding strategy.
\end{enumerate}

\textbf{Q5. Explain concept of crossover and methods possible in GA.}
\begin{enumerate}
\item Crossover combines two or more parent chromosomes.
\item Purpose is to create better offspring by recombination.
\item It is applied after selection stage.
\item Single-point crossover swaps segments at one cut point.
\item Two-point crossover swaps middle segment between two cuts.
\item Multi-point crossover swaps multiple alternating segments.
\item Uniform crossover performs locus-wise random mixing.
\item Three-parent crossover involves additional parent guidance.
\item Shuffle crossover randomizes order before recombination.
\item Precedence-preservative crossover protects order constraints.
\item Ordered and partially matched crossover are used for permutations.
\item Crossover rate controls how often recombination is applied.
\item Excessive crossover without diversity control can destabilize search.
\item Conclusion: crossover is the primary exploration mechanism in GA.
\end{enumerate}

\textbf{Q6. Apply Mamdani FIS to design washing-machine wash-time controller (input: dirt, grease).}
\begin{enumerate}
\item Define input linguistic variables: dirt and grease.
\item Choose three descriptors each (low, medium, high).
\item Define output variable wash-time with five descriptors (very short to very long).
\item Assign membership functions for all input and output labels.
\item Construct rule base, e.g., IF dirt high AND grease high THEN wash-time very long.
\item Add all rule combinations to cover full input space.
\item Fuzzify measured crisp dirt and grease values.
\item Compute rule strengths using AND operator (typically min/product).
\item Clip/scaled output sets using rule strengths.
\item Aggregate all clipped output sets (typically max).
\item Obtain one combined fuzzy wash-time output distribution.
\item Apply defuzzification (centroid/weighted method).
\item Resulting crisp output gives final wash-time decision.
\item Conclusion: Mamdani controller converts expert washing rules into practical control output.
\end{enumerate}

\section*{Formula Sheet (Module 4)}
\begin{itemize}
\item FIS rule: IF $x$ is $A$ AND $y$ is $B$ THEN $z$ is $C$ (Mamdani)
\item Sugeno rule: IF $x$ is $A$ AND $y$ is $B$ THEN $z=ax+by+c$
\item Roulette probability: $P_i=f_i/\sum_j f_j$
\item Common GA flow: init $\to$ fitness $\to$ selection $\to$ crossover $\to$ mutation $\to$ loop
\item Stopping check: max generation, target fitness, stagnation, time limit
\end{itemize}

\section*{One-Page Quick Revision Summary}
\textbf{Memory hooks}
\begin{itemize}
\item FIS = fuzzy decision engine with 5 classic blocks.
\item Mamdani = fuzzy consequents, Sugeno = function consequents.
\item GA operators order: selection $\to$ crossover $\to$ mutation.
\item Encoding and crossover must be compatible.
\item Stopping criteria should be stated explicitly in every GA long answer.
\end{itemize}

\textbf{Exam rapid skeleton}
\begin{enumerate}
\item One-line definition.
\item Block diagram.
\item Steps/algorithm list.
\item One example rule or chromosome example.
\item Final application statement.
\end{enumerate}

\examtip{For diagram-heavy questions, draw compact block diagrams first, then explain each block in 1 line; this secures structure marks fast.}

\end{document}
